We compare three ways of identifying a dynamics model from noisy state measurements for control: sequentially thresholded sparse regression (SINDy-STLSQ), a small MLP neural ODE, and linear least squares with control (DMDc), all reimplemented here, with the true model as a reference. On an inverted damped pendulum and the Van der Pol oscillator we vary measurement noise and data quantity and report open-loop multi-step prediction error and the cost of LQR and sampling-based MPC controllers designed on each identified model and run on the true system. In this small CPU study (5 seeds, one fixed hyperparameter setting per method, two one-dimensional sweeps) SINDy has the lowest prediction error of the three identified models at all noise levels on the pendulum and up to noise 0.05 on Van der Pol with 1000 transitions. SINDy also fails: on Van der Pol its prediction diverges in one of five seeds at noise 0.1 and in two of five seeds with 200 transitions, where it beats the MLP in only 3 of 5 seeds, so the small-data part of our first hypothesis is refuted on Van der Pol by the criterion registered in advance. The absolute gap to the MLP widens with noise up to 0.05 on both systems (up to 0.1 on Van der Pol) instead of shrinking as hypothesised. On Van der Pol closed-loop cost separates the models far less than prediction error: all models, including the linear one, are within 6.7\% of the true-model LQR cost. On the pendulum the linear model fails and the neural model degrades with noise and is very poor with 200 transitions. Dropping the trigonometric library terms degrades SINDy badly on the pendulum, smoothing helps on the pendulum but hurts on Van der Pol, and the sparsity threshold matters little between 0.01 and 0.3. Results are limited to two low-dimensional systems, state measurement only, and untuned baselines.
